\section{Trie}
\whenuse{Mantener un diccionario de cadenas y consultar palabras exactas o
cuántas comienzan con cierto prefijo en $O(|s|)$. Útil para autocompletado,
prefijos compartidos y XOR máximo mediante un Trie binario. Conviene con
alfabetos pequeños; puede consumir mucha memoria si hay pocas coincidencias
entre prefijos.}
Versi\'on para letras min\'usculas inglesas.

\begin{lstlisting}[style=cpp]
struct Trie {
    struct Node {
        array<int, 26> next;
        int terminal = 0, passing = 0;
        Node() { next.fill(-1); }
    };
    vector<Node> nodes{Node()};

    void insert(const string &word) {
        int v = 0;
        ++nodes[v].passing;
        for (char c : word) {
            int letter = c - 'a';
            assert(0 <= letter && letter < 26);
            if (nodes[v].next[letter] == -1) {
                nodes[v].next[letter] = nodes.size();
                nodes.emplace_back();
            }
            v = nodes[v].next[letter];
            ++nodes[v].passing;
        }
        ++nodes[v].terminal;
    }

    int count_word(const string &word) const {
        int v = 0;
        for (char c : word) {
            int letter = c - 'a';
            if (letter < 0 || letter >= 26 || nodes[v].next[letter] == -1)
                return 0;
            v = nodes[v].next[letter];
        }
        return nodes[v].terminal;
    }
};
\end{lstlisting}

\complexity{$O(|s|)$ por inserci\'on o b\'usqueda; memoria proporcional a la suma
de longitudes, hasta $O(\Sigma|s|)$.}

Cada nodo representa un prefijo. \texttt{terminal} cuenta cuántas veces se
insertó exactamente la palabra que termina allí; \texttt{passing} cuenta
cuántas palabras, incluyendo repetidas, pasan por ese prefijo.

Ejemplo de uso:

\begin{lstlisting}[style=cpp]
Trie trie;
trie.insert("car");
trie.insert("cat");
trie.insert("car");
trie.insert("dog");
cout<<trie.count_word("car")<<'\n';   // 2
cout<<trie.count_word("ca")<<'\n';    // 0
cout<<trie.count_prefix("ca")<<'\n';  // 3
cout<<trie.count_prefix("do")<<'\n';  // 1
cout<<trie.count_prefix("x")<<'\n';   // 0
\end{lstlisting}

\texttt{count\_word(s)} exige que la palabra completa esté insertada;
\texttt{count\_prefix(s)} cuenta todas las palabras que comienzan con
\texttt{s}.